In Modena (Table~\ref{tab:singlemodena}), the three-hour single classifier reaches pooled F1 of 0.8509, compared with 0.8669 for the hybrid. Its drift and noise means are 0.7627 and 0.6742; the hybrid reaches 0.8274 and 0.8439. Replay also improves from 0.9060 to 0.9283 in the hybrid. These gains raise family macro F1 from 0.8593 to 0.9045 and retain every hybrid family mean above 0.80. The simpler classifier has its own advantages: higher random and targeted means, and a clean FPR of 0.0317\%, compared with 0.0584\% for the hybrid. The benefit of the additional structure is stronger temporal-family coverage and pooled performance, with a higher clean-alarm rate at these operating points.

Family-balanced calibration exposes the single classifier's threshold trade-off in Modena. At three hours it raises drift and noise to 0.8311 and 0.8141, but replay falls to 0.7772 and pooled F1 to 0.7787; clean FPR rises to 0.1921\%. Thus the two prespecified single-threshold objectives recover different parts of the family profile. The hybrid preserves replay and pooled performance while retaining strong drift/noise means. Table~\ref{tab:balancedbaseline} supplies all results for the alternative calibration objective.
